% Original analytic teaching diagram, 2026-11-12.
% One die throw, two records: X(w)=w (top) and the parity indicator I (bottom).
\definecolor{VB}{HTML}{7188A3}\definecolor{VBF}{HTML}{E5EAF0}
\definecolor{VG}{HTML}{829C8B}\definecolor{VGF}{HTML}{E8EEE9}
\definecolor{VY}{HTML}{B4A16F}\definecolor{VYF}{HTML}{F1ECDD}
\definecolor{VR}{HTML}{A7787D}\definecolor{VRF}{HTML}{F0E5E6}
\definecolor{VP}{HTML}{9588A5}\definecolor{VPF}{HTML}{ECE8F0}
\definecolor{VI}{HTML}{3B4252}\definecolor{VM}{HTML}{6B7280}\definecolor{VL}{HTML}{4C4D4F}
\begin{tikzpicture}[x=1mm,y=1mm,>=Stealth,font=\figurefont\fontsize{8}{10}\selectfont,
lab/.style={text=VM,font=\figurefont\fontsize{7.5}{9}\selectfont},
om/.style={draw=VL,fill=white,rounded corners=1pt,minimum size=6mm,inner sep=0pt,text=VI},
arr/.style={->,VL,line width=.6pt,shorten >=1.2pt,shorten <=1.2pt}]
% number line of X
\draw[VL,line width=.6pt] (8,44) -- (92,44);
\foreach \k in {1,...,6} {
  \fill[VB] ({15+14*(\k-1)},44) circle (1.1);
  \node[above,text=VI] at ({15+14*(\k-1)},45.2) {$\k$};
}
% outcomes
\foreach \k in {1,...,6} {
  \node[om] (w\k) at ({15+14*(\k-1)},24) {$\k$};
  \draw[arr] (w\k.north) -- ({15+14*(\k-1)},42.8);
}
% indicator values
\draw[VL,line width=.6pt] (8,4) -- (92,4);
\fill[VY] (36,4) circle (1.3); \node[below,text=VI] at (36,2.6) {$0$};
\fill[VY] (64,4) circle (1.3); \node[below,text=VI] at (64,2.6) {$1$};
\foreach \k in {1,3,5} { \draw[arr] (w\k.south) -- (36,5.4); }
\foreach \k in {2,4,6} { \draw[arr] (w\k.south) -- (64,5.4); }
\node[anchor=east,text=VI] at (6,44) {$X$};
\node[anchor=east,text=VI] at (6,24) {$\Omega$};
\node[anchor=east,text=VI] at (6,4) {$I$};
\end{tikzpicture}
